\uuid{RvdV}
\titre{Espaces supplémentaires/Noyau/Image}
\auteur{hanani ; mzari}
\organisation{exo7}
\licence{CC BY-NC-SA 4.0}
\datecreate{2020-01-20}
\niveau{L1}
\module{Algèbre}
\chapitre{Application linéaire}
\sousChapitre{Définition}
\difficulte{3}
\theme{Noyau et image, Espaces supplémentaires, Noyau, Image}
\exo7qcmid{lille-2.applications-lineaires.028}

\contenu{
\begin{question}
\qclassification{107.01}
\qtags{motcle=Espaces supplémentaires/Noyau/Image}
Soit $E$ un espace vectoriel, $F$ et $G$ deux sous-espaces supplémentaires dans $E$ et $f$ l'application de $E$ dans $E$ définie par : $$\begin{array}{rccc}f:&E=F\oplus G&\to&E\\
& x=x_1+x_2, \; (x_1\in F, x_2 \in G)&\to &x_1.  \end{array}$$
$f$ est appelée la projection vectorielle de $E$ sur $F$ parallèlement à $G$. Quelles sont les assertions vraies ?
\begin{answers}
\good{$f$ est un endomorphisme de $E$}.
\bad{$f^2=0$}.
\good{$f^2=f$}.
\good{$F=\Im f$ et $G=\ker f$}.
\end{answers}
\begin{explanations} On vérifie que $f$ est un endomorphisme, $f^2=f$, $\ker f=G$ et $\Im f=F$.
\end{explanations}
\end{question}
}
